\documentclass{article}
\usepackage{decomp}
\begin{document}

\decompnote{Full verbatim transcript, produced tool-first: the text was placed from a Datalab API (Chandra) conversion of the scan made on 2026-10-18 (\texttt{work/datalab/glauert1958.md}) by the script \texttt{work/build\_main.py}, and proofread against the page images by the decompiler; every correction is in the script's correction table and in the README. The source is a 2003 scan of the printed journal pages (5 pages: one 300-ppi bilevel image per page, plus a 150-ppi half-tone of the plate) with an OCR text layer; the page images are authoritative. No licence is printed in the PDF. Running headers and page numbers (191--194) are left out; the journal line printed at the foot of p.\,191 is given once below. Footnotes are attached at their marks.}

\sourcefile{glauert1958.pdf}

\pdfpage{1}
\noindent\textsc{J. Biophysic. and Biochem. Cytol., 1958, Vol. 4, No. 2}\par
\decompnote{The journal line above is printed at the foot of the first page (p.\,191).}

\section*{Araldite as an Embedding Medium for Electron Microscopy}

\noindent\textsc{By Audrey M. Glauert* and R. H. Glauert, Ph.D.}

\noindent\textit{(From the Strangeways Research Laboratory and the University Chemical Laboratory, Cambridge, England)}

\begin{center}
\textsc{Plate 98}

(Received for publication, December 10, 1957)
\end{center}

\noindent * Sir Halley Stewart Research Fellow.
\decompnote{The asterisk note is printed as a footnote at the foot of the left column of p.\,191; placed here after the author line.}

\section*{Abstract}

Epoxy resins are suitable media for embedding for electron microscopy, as they set uniformly with virtually no shrinkage. A mixture of araldite epoxy resins has been developed which is soluble in ethanol, and which yields a block of the required hardness for thin sectioning. The critical modifications to the conventional mixtures are the choice of a plasticized resin in conjunction with an aliphatic anhydride as the hardener. The hardness of the final block can be varied by incorporating additional plasticizer, and the rate of setting can be controlled by the use of an amine accelerator. The properties of the araldite mixture can be varied quite widely by adjusting the proportions of the various constituents. The procedure for embedding biological specimens is similar to that employed with methacrylates, although longer soaking times are recommended to ensure the complete penetration of the more viscous epoxy resin. An improvement in the preservation of the fine structure of a variety of specimens has already been reported, and a typical electron microgram illustrates the present paper.

\section*{Introduction}

Since the first description of a mixture of araldite epoxy resins as an embedding medium for electron microscopy (Glauert, Rogers, and Glauert, 1956), it has come into routine use in a number of laboratories (Huxley, 1957; Birbeck and Mercer, 1957). Only a short outline of the method of embedding was given in the original communication, and it is the purpose of this paper to describe further details of the technique and to indicate some of the physical and chemical properties of the mixture.

\subsection*{Epoxy Resins}

Although the first epoxy resins were made in the late 1930's, the commercial development of these polymers has depended on the ready availability of certain raw materials and dates only from the end of the Second World War.

Chemically the epoxy resins are polyaryl ethers of glycerol bearing terminal epoxy groups and may be given the general formula (I),

\begin{center}
\includegraphics[width=0.6\linewidth,keepaspectratio]{figures/formula1.png}
\end{center}
\alttext{Structural formula (I), p.\,191, drawn in text characters on two lines with the label ``(I)'' at the right. First line: an epoxide ring (the letter O above an inverted V of two bonds) sits over the first two carbons of ``CH\textsubscript{2}CHCH\textsubscript{2}O''; a bond then leads into a square bracket enclosing ``ROCH\textsubscript{2}CHCH\textsubscript{2}O'' in which an ``OH'' is bonded above the middle CH; the closing bracket carries the subscript \emph{n}, and a bond leads out of it. Second line (set lower right): ``ROCH\textsubscript{2}CHCH\textsubscript{2}'' with an epoxide O bridging its last CH and CH\textsubscript{2}. Read atom by atom: a 2,3-epoxypropyl (glycidyl) ether at each end, joined through \emph{n} repeats of --R--O--CH\textsubscript{2}--CH(OH)--CH\textsubscript{2}--O-- and a final --R--O--, i.e. epoxide--CH\textsubscript{2}--O--[R--O--CH\textsubscript{2}--CH(OH)--CH\textsubscript{2}--O--]\emph{n}--R--O--CH\textsubscript{2}--epoxide, with R the divalent group (II). SMILES with R = (II): for \emph{n} = 0, C1OC1COc1ccc(cc1)C(C)(C)c1ccc(cc1)OCC1CO1; for \emph{n} = 1, C1OC1COc1ccc(cc1)C(C)(C)c1ccc(cc1)OCC(O)COc1ccc(cc1)C(C)(C)c1ccc(cc1)OCC1CO1; each further repeat inserts another OCC(O)COc1ccc(cc1)C(C)(C)c1ccc(cc1) unit. No charges; all bonds single except the aromatic rings of R. The subscript digits of the scan are small and the ``n'' is italic.}

in which R is commonly the diphenyl propane system (II) and \textit{n} is an integer.

\begin{center}
\includegraphics[width=0.5\linewidth,keepaspectratio]{figures/formula2.png}
\end{center}
\alttext{Structural formula (II), p.\,191, with the label ``(II)'' at the right: two benzene rings drawn as hexagons with alternating double bonds, each para-disubstituted, joined through a central carbon ``C'' that carries ``CH\textsubscript{3}'' above and ``CH\textsubscript{3}'' below it (the four bonds of C: left ring, right ring, upper CH\textsubscript{3}, lower CH\textsubscript{3}). A free bond leaves each ring at its outer para position towards the left and right edges of the drawing: these are the two attachment points of R in formula (I). SMILES of the divalent unit, attachment points as asterisks: *c1ccc(cc1)C(C)(C)c1ccc(*)cc1. No charges.}

The resins range from viscous liquids to fusible solids depending on the molecular weight, \textit{i.e.} the value of \textit{n}. They may be cured by a variety of bifunctional setting agents which add across the epoxy groups of the resin molecules to give three-dimensional structures. For cold setting these setting agents are commonly aliphatic polyamines, and for hot setting aromatic anhydrides, but the number of possible setting agents is very large and this enables a wide range of products with very different characteristics to be obtained. Significantly, the setting can be made to occur with a volume shrinkage as low as 2 per cent, a comparable figure for methacrylate being 15 to 20 per cent. This relatively small change in volume is due, on the one hand, to the setting step being an addition process and, on the other, to the fact that the resin is highly associated in the uncured state.

These materials are commonly used as surface coating agents, as casting and encapsulating materials, as adhesives, and for many other purposes.

\pdfpage{2}
\subsection*{Epoxy Resins as Embedding Media for Electron Microscopy}

The facts that the epoxy resins set uniformly and with substantially no change of volume indicate that they should be suitable as embedding media for electron microscopy. This was realised by a number of different workers who investigated the standard products of various firms. The most important contribution was made by Maaløe and Birch-Andersen (1956), who reported their results at a symposium of the Society for General Microbiology in London in 1956.

Maaløe and Birch-Andersen embedded their specimens in a standard liquid epoxy resin, which they designated EPO, in conjunction with an aliphatic polyamine (diethylene triamine) as hardener. With this mixture they obtained sections of \textit{Salmonella typhimurium} in which the cells were free of the damage and distortion that so often occurs on embedding in methacrylate.

The EPO mixture is very viscous and has only a limited compatability with ethanol. Consequently, it is difficult to handle and a prolonged embedding procedure is necessary to ensure that the resin penetrates fully into the specimen. Maaløe and Birch-Andersen found that EPO often failed to diffuse into the tubular ``nuclear'' structures of \textit{Salmonella typhimurium}.

\subsection*{The Araldite Mixture}

In order to avoid some of the difficulties encountered by Maaløe and Birch-Andersen (1956), a series of different mixtures based on araldite\footnote{Trade name for epoxy resins developed by Messrs. Ciba, Ltd., Basle, and made in Great Britain by Aero Research, Ltd., Duxford, Cambridge.} epoxy resins was investigated. It was found that a mixture incorporating an aliphatic anhydride as the hardener had many desirable properties. The hardener, being of high molecular weight, needed to be mixed with the resin in almost equal amounts to give the correct stoichiometrical proportion. This resulted in a mixture with a much reduced viscosity compared with one incorporating an aliphatic polyamine, in which the hardener addition is only about 10 per cent. This mixture proved also to be freely soluble in ethanol.

Although the reaction between the anhydride and the epoxy resin normally requires a high temperature for curing, the addition of an amine accelerator enables this temperature to be reduced to a value suitable for the embedding process. The hardness of the final block can be controlled by the addition of a suitable amount of dibutyl phthalate which is a plasticizer for the epoxy resin.

As a result of these trials with araldite resins, the following mixture was recommended as an embedding medium:—

\begin{center}
\begin{tabular}{@{}lr@{}}
Araldite M & 10.0 ml.\\
Hardener 964 B & 10.0 ml.\\
Dibutyl phthalate & 1.0 ml.\\
Accelerator 964 C & 0.5 ml.\\
\end{tabular}
\end{center}

This mixture sets in about 48 hours at 48°C. and yields a light gold block that has a similar hardness to methacrylate (methyl/\textit{n}-butyl:15/85). The hardness can be varied throughout the methacrylate range by adjusting the amount of dibutyl phthalate that is added. Care has to be taken, however, not to add too much dibutyl phthalate, or the addition reaction will be incomplete and the final block will be unstable.\footnote{\textit{I.e.}, the block will not be sufficiently rigid for cutting and may suffer dimensional changes with time.}

The time of setting is critically dependent on the amount of accelerator that is present, so that it is necessary to measure this amount accurately. An excess of the accelerator can alter the characteristics of the final block which becomes dark in colour and too brittle for satisfactory cutting.

\subsection*{Embedding in Araldite}

(a) \textit{Biological Specimens.}—The araldite mixture is readily soluble in ethanol, so that the same fixation and dehydration procedure can be followed as for embedding in methacrylate. Normally the specimen is passed from pure ethanol to a 50/50 mixture of ethanol and araldite, but it is also possible to pass directly from 90 per cent ethanol to the araldite mixture, if this is thought to be an advantage (D. Lacy, personal communication).

The viscosity of araldite is somewhat greater than that of methacrylate, but it can be reduced by raising the temperature. It is recommended that the initial soakings and final embedding of the specimen should be carried out at a stable temperature of 48°C. At higher temperatures the lowered viscosity of the mixture is offset by the increased rate of setting. Also there is a tendency for soft specimens to rise in the capsule during the setting process at temperatures above 50°C. If it is thought necessary to avoid heating the specimen to 48°C., the soaking can be done at a \pdfpage{3}lower temperature, but then the times of soaking must be increased considerably to ensure that the resin penetrates completely into the specimen. It is advisable to make up the mixture without the accelerator for the preliminary soakings, and then there is no danger of the mixture starting to set too soon.

The following table indicates some typical soaking times that have been found to be adequate.

\begin{center}
\begin{tabular}{@{}lcc@{}}
\toprule
 & \begin{tabular}[t]{@{}c@{}}Soft specimens,\\ \emph{e.g.} lung, liver,\\ spleen, etc.\end{tabular} & \begin{tabular}[t]{@{}c@{}}Harder specimens:\\ muscle, bacteria,\\ etc.\end{tabular}\\
\midrule
50/50 ethanol/araldite & 1--2 hrs. at 48°C. & 4--6 hrs. at 48°C.\\
Araldite, less accelerator & 2--3 hrs. at 48°C. & 24 hrs. at 48°C.\\
Araldite, with accelerator & 2--3 hrs. at 48°C. & 24 hrs. at room temperature\\
\bottomrule
\end{tabular}
\end{center}

These times are merely given as an indication of the range of soaking periods that has been tried. The scheme can obviously be modified to meet the requirements of many different types of specimen.

(b) \textit{Other Specimens.}—Little work has yet been done on the study of sections of very hard biological and non-biological specimens embedded in epoxy resins. Glauert, Rogers, and Glauert (1956) were able to cut sections of adult human hair embedded in araldite. The embedding medium is not expected to penetrate such specimens, so that the higher viscosity of araldite is not a disadvantage. Also the fact that the mixture does not shrink on setting indicates that it will ``hold'' hard materials far more firmly than methacrylate, which tends to retract away from the surface. A combination of embedding in araldite and sectioning with a diamond knife should be excellent for metal specimens.

\subsection*{Technical Details of Embedding in Araldite}

Epoxy resins are strongly adhesive materials and should be handled with some care. Although the uncured resin can be removed by washing with ethanol, after curing it is completely insoluble. In these circumstances it is advisable to keep a separate set of glassware for measuring and mixing the resins.

Araldite M, Hardener 964 B, and dibutyl phthalate will keep for many months at room temperature. It is best to store the resins in glass-stoppered bottles and to wipe the stoppers carefully each time some resin has been poured out. Accelerator 964 C may deteriorate with time and should be kept as dry as possible.

For ease of handling it is convenient to heat the resin and its hardener, say to 60°C., before they are mixed. The mixture is most easily prepared by pouring each component in turn into a warm graduated cylinder. It can then be poured into a small conical flask and stirred well with a glass rod. Some air bubbles develop as the mixture is stirred, but these will go if the mixture is allowed to stand for a few minutes. It is particularly important that the final mixture that is poured into the capsules should be free of bubbles. The mixture without the accelerator will keep for some weeks at 0°C.

If gelatin capsules are used for the final embedding, they should be clean and dry. It is convenient to remove the specimen from the last soaking mixture with a fine glass pipette and place it at the exact bottom of a capsule. The capsule can then be filled immediately with the araldite mixture (with accelerator). Sometimes the specimen is pushed to one side as the liquid is poured in, but it can be guided back to the centre with a fine glass pipette before the top is placed on the capsule. Care must again be taken not to introduce any bubbles of air at this stage.

\subsection*{Sectioning of Araldite}

The blocks of araldite are light gold in colour and show none of the shrinkage effects associated with methacrylate. Thin sections have been cut with glass and diamond knives on a variety of microtomes with essentially the same technique as that employed for obtaining sections of methacrylate blocks. The hardened araldite is not soluble in ethanol or acetone, and the sections are routinely cut onto a pure water surface.

\subsection*{Electron Microscopy of Araldite Sections}

Unlike methacrylate, epoxy resins are not degraded under electron bombardment and araldite sections do not ``clear'' in the electron beam. Consequently the araldite will not add to the contamination within the microscope. Araldite also appears to have a somewhat greater electron density than methacrylate, but sufficient contrast can be obtained if care is taken to make the sections thin enough.

\section*{Conclusions}

Sufficient work has now been done in a number of different laboratories to indicate that araldite is an attractive alternative for methacrylate as an embedding medium for electron microscopy. A \pdfpage{4}suitable choice of resin, hardener, and accelerator has resulted in a mixture which has many of the required properties of an embedding medium and which is quite easy to handle during routine use. An improvement in the preservation of the fine structure of a variety of specimens has already been reported (Glauert and Brieger, 1956; Huxley, 1957; Birbeck and Mercer, 1957), and a typical electron micrograph is given here. There seems no reason to doubt that similar successes will be achieved when araldite is applied in other fields of investigation.

\textit{Sources of Supply.}—Great Britain, Aero Research Ltd., Duxford, Cambridge; Europe, Ciba Ltd., Basle, Switzerland; United States of America, through Ciba Co., Inc., Plastics Division, Kimberton, Pennsylvania.

\section*{References}

\begin{hangpars}
Birbeck, M. S. C., and Mercer, E. H., Some applications of an epoxide embedding medium, Demonstration at 15th Meeting of Electron Microscope Society of America, Boston, 1957.

Glauert, A. M., and Brieger, E. M., Ultra-thin sections of avian tubercle bacilli in a new embedding medium, Proceedings Stockholm Conference on Electron Microscopy, Stockholm, Almqvist \& Wiksell, 1956, 111.

Glauert, A. M., Rogers, G. E., and Glauert, R. H., A new embedding medium for electron microscopy, \textit{Nature}, 1956, \textbf{178}, 803.

Huxley, H. E., Some observations on the use of Araldite as an embedding medium, Paper at 15th Meeting of Electron Microscope Society of America, Boston, 1957.

Maaløe, O., and Birch-Andersen, A., On the organization of the ``nuclear material'' in \textit{Salmonella typhimurium}, 6th Symposium of the Society for General Microbiology, Bacterial Anatomy, Cambridge, The University Press, 1956, 261.
\end{hangpars}

\section*{Explanation of Plate 98}

\textbf{Fig. 1.} Section through a group of bovine tubercle bacilli embedded in araldite. The outer cell wall has remained closely apposed to the cytoplasmic membrane, and fine cytoplasmic structures have been preserved. The large dense granules at the ends of the cells are metachromatic, metaphosphate bodies and appear vacuolated as a result of bombardment with the electron beam. (From an unpublished study by E. M. Brieger and Audrey M. Glauert). × 100,000.

\pdfpage{5}
\noindent\textsc{The Journal of Biophysical and Biochemical Cytology}\hfill\textsc{Plate 98}\\
\mbox{}\hfill\textsc{Vol. 4}
\decompnote{Plate page (PDF page 5): the lines above are its header; the figure number 1 is printed inside the half-tone (lower left corner), and the line after the description is the plate's footer. The caption (Explanation of Plate 98) is printed on p.\,194, before the plate, and is kept there.}
\begin{center}
\includegraphics[width=0.8\linewidth,height=0.75\textheight,keepaspectratio]{figures/fig1.png}
\end{center}
\alttext{Transmission electron micrograph, greyscale, portrait format (the embedded half-tone is 805 by 987 pixels at 150 ppi), the figure number ``1'' printed in the lower left corner; no scale bar (the explanation gives \texttimes\ 100,000). Thin section of bacteria: the background (embedding medium) is light grey. Two large rod-shaped cell profiles cross the middle of the field horizontally, one above the other, nearly touching along their facing edges; each is rounded at its right end and runs into neighbouring cells or out of the field at the left. Each rod is bounded by a smooth, continuous outer contour with a narrow lighter outer zone (the cell wall) that follows the outline of the darker cytoplasm closely, with no gap between them; a separate membrane line is hard to make out at this reproduction. The cytoplasm is medium grey with a mottled, finely granular texture and some darker patches: in the upper rod a darker region in its left third, in the lower rod a dark patch left of centre and a small ring-shaped light spot right of centre. Near the right end of each rod lies one large, round, very dark granule, roughly a third to a half of the cell width across, peppered with small light round spots (the ``vacuolated'' look of the explanation); the upper rod's granule is the larger and nearly circular, the lower rod's is smaller and kidney-shaped. Other profiles at the edges: top right, part of a rounded cell with a solid, uniformly dark round granule (without light spots) and a lighter rim; upper left and left middle, two faint pale rounded profiles cut by the frame; bottom left, a rod cut obliquely, running from the lower left corner up to the right, darker and coarser, with a few black specks near its lower edge; bottom right, a rounded cell end with granular cytoplasm and a darker outer zone. No labels, arrows or annotations other than the figure number. Not a plotted panel: no data file.}

\noindent (Glauert and Glauert: Araldite as embedding medium)

\end{document}
